\documentclass[10pt,twocolumn,letterpaper]{article}

\usepackage[utf8]{inputenc}
\usepackage[english]{babel}
\usepackage{amsmath, amssymb, amsfonts}
\usepackage{graphicx}
\usepackage{hyperref}
\usepackage{booktabs}
\usepackage{geometry}
\geometry{a4paper, margin=1.5cm}
\usepackage{caption}
\usepackage{subcaption}
\usepackage{xcolor}

\title{\Large \bf Empirical Analysis of Learning Dynamics: \\ Double Descent, Neural Collapse and Internal Representations via CKA}
\author{Marco Bardazzi}
\date{\today}

\begin{document}

\maketitle

\begin{abstract}
This study explores the complex optimization and generalization dynamics of deep neural networks through the analysis of a one-dimensional dataset (MNIST-1D) corrupted with 20\% label noise. We investigate the presence of the \textit{Deep Double Descent} phenomenon as the model width varies, demonstrating how \textit{Neural Collapse} metrics (NC1-NC4) faithfully reflect the test error. Furthermore, we evaluate the theoretical limit of Neural Collapse as feature dimensionality varies and use \textit{Centered Kernel Alignment} (CKA) to measure representation similarity across different architectures, training epochs, and initializations (seeds), revealing that deep layers undergo greater drift and exhibit high geometric degeneracy.
\end{abstract}

\section{Introduction}
Classical machine learning theory suggests a "U"-shaped bias-variance trade-off, where overfitting monotonically worsens performance on unseen data as model capacity increases. However, modern neural networks exhibit a \textit{Double Descent} behavior \cite{belkin2019}: past the interpolation threshold, heavily over-parameterized networks generalize excellently. 

In parallel, Papyan et al. \cite{papyan2020} demonstrated that in the terminal phase of training, a geometric phenomenon known as \textit{Neural Collapse} (NC) emerges, characterized by a contraction of within-class variance (NC1) and the convergence of class means to a \textit{Simplex Equiangular Tight Frame} (ETF) (NC2). 

This study replicates and analyzes these phenomena using variants of CNN and ResNet1D, enriching the investigation with the study of NC metrics as width varies and the comparison of geometric similarity via CKA \cite{kornblith2019}.

\section{Methodology}

\subsection{Dataset and Noise}
The experiment uses MNIST-1D, a low-dimensional synthetic dataset that retains the morphological properties of deep learning. To induce the Double Descent phenomenon, a \textbf{20\% label noise} was introduced. The noise forces the model to exhaust its capacity to memorize incorrect labels in the critical regime, exacerbating the generalization error peak.

\subsection{Architectures and Optimization}
Two families of models were compared: \texttt{ResNet1D} (with residual connections and BatchNorm) and \texttt{StandardCNN1D}. The experiments were conducted by varying the width parameter $W \in [2, 64]$. Two optimizers were compared: SGD with momentum (0.9) and Cosine Annealing, and Adam.

\subsection{Evaluation Metrics}
In addition to the Test Error, the 4 Neural Collapse metrics were calculated:
\begin{itemize}
    \item \textbf{NC1 (Variability):} Ratio between within-class and between-class covariance.
    \item \textbf{NC2 (Simplex ETF):} Frobenius distance between the covariance matrix of class means and the ideal Simplex ETF.
    \item \textbf{NC3 (Self-Duality):} Alignment between classifier weights and class feature means.
    \item \textbf{NC4 (NCC Mismatch):} Percentage disagreement between network predictions and a \textit{Nearest Class Center} classifier.
\end{itemize}
For the analysis of internal representations, \textbf{Linear CKA} was used, which evaluates the similarity between two activation matrices by normalizing their cross-covariance with respect to their individual variances.

\begin{figure*}[h!]
    \centering
    \includegraphics[width=1.0\textwidth]{neural_collapse_cka_results.png}
    \caption{Complete experimental results. Row 1 shows the evolution of the error and Neural Collapse (NC1-NC3) as parameters vary (Exp 1). Row 2 illustrates NC4 (left), \textit{epoch-wise} evolution (center), and generalized collapse with respect to dimensions (right). Row 3 shows layer-by-layer CKA analyses.}
    \label{fig:results}
\end{figure*}

\section{Results and Interpretation}

\subsection{Exp 1: Double Descent and Neural Collapse (Model-wise)}
The graphs in the first row (Fig. \ref{fig:results}) clearly show the Double Descent curve. Around $10^3-10^4$ parameters, the Test Error experiences a sudden peak. This is the \textbf{critical interpolation regime}: the model has exactly the degrees of freedom needed to memorize the 20\% noise, resulting in highly non-linear and unstable functions that fail on test data.
Further increasing the parameters ($>10^4$), the "blessing of dimensionality" takes over: the inductive bias of the optimizer (particularly SGD) favors solutions with smaller norms and smoother decision boundaries, reducing the error.

The most relevant result is the \textbf{perfect correlation between Test Error and the NC1, NC2, NC3, and NC4 metrics}. At the critical point, Neural Collapse "breaks": the high effort to memorize noise destroys the geometric symmetry of the latent space. Past the critical threshold, representations collapse again towards a perfect Simplex ETF. It is also observed that \textbf{SGD heavily dominates Adam} in highly over-parameterized regimes, favoring a more extreme Neural Collapse.

\subsection{Exp 2: Epoch-wise Dynamics}
Analyzing the error across epochs (Fig. \ref{fig:results}, row 2, center), the "Critical" model ($W=12$, located exactly on the Double Descent peak) shows massive fluctuations. Training on a noisy dataset with borderline capacity renders the loss surface fragmented. Conversely, the "Large" model ($W=64$) quickly learns robust features in early epochs and memorizes noise only later, maintaining a low and stable test error.

\subsection{Exp 3: Dimensional Limits of NC}
Geometric theory dictates that, for $K$ classes, a perfect Simplex ETF can only exist if the feature dimension $d \ge K-1$. Since MNIST-1D consists of 10 classes, we expect collapse to occur for $d \ge 9$.
The "Gen NC" graphs confirm this theory: for $d < 9$, NC2 is high because the latent space lacks sufficient dimensions to accommodate 10 symmetric equiangular vectors. Exactly at the red dashed line ($K-1=9$), the NC2 metric drops vertically towards zero, empirically confirming the spatial limits predicted by the mathematics of Neural Collapse.

\subsection{Exp 4: Representations Analysis (CKA)}
CKA results provide a magnifying glass on the behavior of intermediate layers:

\paragraph{4A. Cross-Architecture (Small vs Large):} 
The heatmap reveals that early layers (L0, L1) have a relatively high CKA, indicating they extract similar low-level features (e.g., edges and 1D signal frequencies). However, moving towards deep layers (L3), CKA drops to $0.34$. This difference is plausible: the large model possesses the degrees of freedom to collapse features towards an ETF, whereas the small model fails to separate classes into compact clusters, leading to completely dissimilar topologies before the linear classifier.

\paragraph{4B. Epoch-wise Evolution (Drift):}
Comparing activations during training against activations at epoch 1, we notice that layer L0 (Stem) remains almost unchanged ($CKA \approx 1.0$). This indicates that basic filters converge rapidly or undergo irrelevant modifications. Conversely, L2 and L3 show strong "drift". Deep representations continue to reorganize over tens of thousands of epochs, driven by the need to minimize NC metrics (compressing within-class variance and separating class centers).

\paragraph{4C. Cross-Seed Reproducibility:}
Training the same model ($W=32$) with different random seeds, a clear hierarchy emerges. L0 is almost deterministic ($CKA = 0.97$). Deep layers (L2, L3) show moderate values ($CKA \approx 0.59$). Being highly over-parameterized, the network possesses a \textbf{geometric degeneracy}: there are infinite ways to rotate a latent space while still forming a perfect Simplex ETF. Since CKA captures structural similarities, different initializations lead to distinct isomorphic rotations of the final space, solving the task with the same accuracy but with dissimilar latent coordinates.

\section{Conclusions}
This empirical study validates and intertwines three of the most fascinating phenomena in modern deep learning. We demonstrated that the inability to generalize at the interpolation point (Double Descent) is intimately linked to the destruction of latent geometry (Neural Collapse). The injection of noise highlighted how massively parameterized networks optimized with SGD manage to recover optimal geometric structures (Simplex ETF) provided that the latent dimension allows it ($d \ge K-1$). Finally, the CKA analysis confirmed that this collapse is a dynamic and degenerative phenomenon, localized almost exclusively in the final layers of the architecture.

\begin{thebibliography}{9}

\bibitem{belkin2019}
M. Belkin, D. Hsu, S. Ma, S. Mandal.
\textit{Reconciling modern machine-learning practice and the classical bias–variance trade-off.}
Proceedings of the National Academy of Sciences, 2019.

\bibitem{papyan2020}
V. Papyan, X. Y. Han, D. L. Donoho.
\textit{Prevalence of neural collapse during the terminal phase of deep learning training.}
Proceedings of the National Academy of Sciences, 2020.

\bibitem{kornblith2019}
S. Kornblith, M. Norouzi, H. Lee, G. Hinton.
\textit{Similarity of neural network representations revisited.}
International Conference on Machine Learning (ICML), 2019.

\end{thebibliography}

\end{document}